\uuid{DrDK}
\titre{Valeur absolue}
\auteur{bodin ; hanani ; mzari}
\organisation{exo7}
\licence{CC BY-NC-SA 4.0}
\datecreate{2020-01-20}
\niveau{L1}
\module{Analyse}
\chapitre{Propriétés de R}
\sousChapitre{Propriétés des nombres réels}
\difficulte{3}
\theme{Propriétés de nombres réels, valeur absolue}
\exo7qcmid{lille-1.reels.017}

\contenu{
\begin{question}
\qclassification{120.03}
\qtags{motcle=valeur absolue}
Quelles sont les assertions vraies, pour tout $x,y\in\Rr$ ?
\begin{answers}
    \bad{$|x-y| \le |x|-|y|$}

    \good{$|x| \le |x-y|+|y|$}

    \bad{$|x+y| \ge |x| + |y|$}

    \good{$|x-y| \le |x| + |y|$}
\end{answers}
\begin{explanations}
L'inégalité triangulaire est $|x+y| \le |x| + |y|$. Les assertions vraies en découlent.
\end{explanations}
\end{question}
}
